\chapter{Generational Handles}
\label{chap:handles}

\section{Motivation}

Decrease-key and arbitrary remove require a way to locate an existing item inside a heap.
If the library searched for the item by pointer identity, targeted operations would become linear and the advertised heap complexity would be misleading.
heapx therefore exposes insertion handles.

The insert-with-handle operation returns a handle.
That handle can later be passed to the decrease-key or remove operations.
Plain \texttt{heapx\_insert()} does not create a public handle.

\section{Handle Representation}

The public handle is a small value type:

\begin{lstlisting}[style=cstyle, caption={Public handle representation.}]
struct heapx_handle {
    uint64_t heap_id;
    size_t slot;
    unsigned generation;
};
\end{lstlisting}

It contains the logical heap identifier, the slot index in the heap's common handle table, and the generation observed when the handle was created.
The handle owns no memory and can be copied by the caller.

\section{Handle Slot Table}

The common base object stores a handle slot table.
Each slot records the caller-owned item pointer, a backend-owned owner pointer, the current generation, a live bit, and a free-list link.

\begin{lstlisting}[style=cstyle, caption={Internal handle slot.}]
struct heapx_handle_slot {
    void *item;
    void *owner;
    unsigned generation;
    int live;
    size_t next_free;
};
\end{lstlisting}

When a handled item is inserted, heapx either reuses a free slot or grows the slot table geometrically.
When the item is removed or extracted, the slot is marked non-live, its generation is incremented, and the slot is pushed onto the free list.

\section{Stale Handle Protection}

Generations prevent stale handles from becoming valid again after slot reuse.
Suppose an item is inserted into slot $s$ with generation $g$.
After the item is removed, the slot generation becomes $g + 1$.
If the same slot is later reused for another item, the old handle still contains generation $g$, so resolution fails.

The heap identifier prevents handles from one heap being used on another heap.
This is tested explicitly in the test suite.

\section{Backend Owner Pointers}

The handle table stores a backend owner pointer.
This pointer is private to the backend and has different meanings depending on the implementation:

\begin{itemize}
    \item the binary heap stores a stable locator object whose \texttt{index} tracks the current array position;
    \item the Fibonacci heap stores the node pointer directly;
    \item the Kaplan heap stores the node pointer directly.
\end{itemize}

Public targeted operations resolve a handle to this owner pointer before dispatching to the backend.
The backend does not need to know about heap identifiers or slot generations; it receives the local object needed to perform the algorithmic repair.

\section{Binary Heap Locators}

The binary heap is array-based, so item positions change during sift-up, sift-down, extract-min, and remove.
A public handle cannot point directly to an array slot because array entries move.
The backend therefore allocates stable locator objects.
Each locator stores the public handle and the current array index.
Whenever an entry is moved, \texttt{binary\_heap\_store\_entry()} refreshes the locator's index.

This is an important implementation detail because it lets the binary heap support the same targeted API as pointer-based heaps.

\section{Failure Behavior}

If handle attachment fails during insertion, the backend rolls back the partially allocated backend metadata and the item is not inserted.
This preserves the public contract: on failed insert, the caller remains the sole owner of the item and the heap is unchanged from the caller's perspective.
